\begin{abstract}
Wireless Body Area Networks (WBANs) operating under the IEEE 802.15.6 standard must balance strict energy efficiency with strong security. Traditional cryptographic key exchange protocols require significant wireless communication, which drains the batteries of implanted medical sensors and increases the risk of exposing sensitive health data. This project evaluates a physiological-signal-based key agreement protocol. Two sensors on the same body extract synchronized heartbeat features, specifically inter-pulse intervals (IPIs), from the ECG signal and generate the cryptographic key locally, so the key is never transmitted over the air. A Bloom filter reconciles the two sensors' feature sets without revealing the raw biometric values, and HMAC tags confirm that both sides hold the same key. Simulation results on the MIT-BIH Arrhythmia Database show that both sensors derive an identical 160-bit key with matching HMAC verification. The Python prototype completes in about 7.04 seconds, most of which is ECG processing. We also validated the Bloom filter experimentally, and found and fixed a defect that had raised its false-positive rate from the 0.096\% target to 49.5\%. The approach removes over-the-air key exchange while keeping the communication overhead to about 100 bytes, which makes it suitable for battery-constrained WBAN devices.

\medskip
\noindent\textbf{Keywords:} WBAN, IEEE 802.15.6, ECG, inter-pulse interval, Bloom filter, key agreement, HMAC.
\end{abstract}
